\documentclass[11pt,a4paper]{article}
\usepackage[utf8]{inputenc}
\usepackage[T1]{fontenc}
\usepackage[french]{babel}
\usepackage[margin=1.8cm]{geometry}
\usepackage{amsmath,amssymb}
\usepackage{xcolor}
\usepackage{colortbl}
\usepackage{tcolorbox}
\usepackage{booktabs}
\usepackage{array}
\usepackage{multicol}
\usepackage{enumitem}
\usepackage{tikz}
\usepackage{pgfplots}
\pgfplotsset{compat=1.17}
\usepackage{fancyhdr}
\usepackage{multirow}
\usepackage{fontawesome5}
\usepackage{longtable}

\tcbuselibrary{skins,breakable}

% ---- Couleurs ----
\definecolor{biocolor}{RGB}{0,110,110}
\definecolor{bioaccent}{RGB}{200,80,40}
\definecolor{biolight}{RGB}{230,245,245}
\definecolor{bioyellow}{RGB}{255,245,210}
\definecolor{corrige}{RGB}{20,120,60}
\definecolor{corrigelight}{RGB}{230,245,235}

% ---- En-têtes ----
\pagestyle{fancy}
\fancyhf{}
\fancyhead[L]{\textcolor{corrige}{\textbf{Corrigé -- Challenge Statistique}}}
\fancyhead[R]{\textcolor{corrige}{STS BIOAL -- Module 5}}
\fancyfoot[C]{\thepage}
\renewcommand{\headrulewidth}{0.6pt}
\renewcommand{\headrule}{\hbox to\headwidth{\color{corrige}\leaders\hrule height \headrulewidth\hfill}}

% ---- Boîtes ----
\newtcolorbox{epreuve}[2][]{
  breakable, colback=corrigelight, colframe=corrige,
  coltitle=white, title={\textbf{#2}},
  fonttitle=\bfseries\large, boxrule=0.8pt, arc=3pt,
  left=6pt, right=6pt, top=6pt, bottom=6pt, #1
}

\newtcolorbox{regle}{
  colback=bioyellow, colframe=bioaccent,
  boxrule=0.8pt, arc=3pt,
  left=6pt, right=6pt, top=6pt, bottom=6pt
}

\newtcolorbox{vignetteCours}[1]{
  breakable, enhanced, colback=biolight, colframe=biocolor,
  coltitle=white, fonttitle=\bfseries, title={\faBookOpen\ \ #1},
  boxrule=0.7pt, arc=4pt, left=6pt, right=6pt, top=5pt, bottom=5pt,
  drop fuzzy shadow=black!15, width=\linewidth
}

\newtcolorbox{vignetteMethode}[1]{
  breakable, enhanced, colback=bioyellow, colframe=bioaccent,
  coltitle=white, fonttitle=\bfseries, title={\faTools\ \ #1},
  boxrule=0.7pt, arc=4pt, left=6pt, right=6pt, top=5pt, bottom=5pt,
  drop fuzzy shadow=black!15, width=\linewidth
}

\newtcolorbox{vignetteAlerte}[1]{
  breakable, enhanced, colback=red!5, colframe=red!70!black,
  coltitle=white, fonttitle=\bfseries, title={\faExclamationTriangle\ \ #1},
  boxrule=0.7pt, arc=4pt, left=6pt, right=6pt, top=5pt, bottom=5pt,
  drop fuzzy shadow=black!15, width=\linewidth
}

\newtcolorbox{reponse}[1][]{
  breakable, enhanced, colback=white, colframe=corrige,
  coltitle=white, fonttitle=\bfseries,
  title={\faCheckCircle\ \ Réponse},
  boxrule=0.7pt, arc=4pt, left=6pt, right=6pt, top=5pt, bottom=5pt, #1
}

% ---- Données ----
\newcommand{\dataS}{%
12{,}1 ; 12{,}6 ; 13{,}0 ; 13{,}3 ; 13{,}6 ; 13{,}9 ; 14{,}0 ; 14{,}2 ;
14{,}4 ; 14{,}5 ; 14{,}7 ; 14{,}9 ; 15{,}0 ; 15{,}1 ; 15{,}2 ; 15{,}3 ;
15{,}4 ; 15{,}5 ; 15{,}6 ; 15{,}8 ; 15{,}9 ; 16{,}0 ; 16{,}1 ; 16{,}3 ;
16{,}4 ; 16{,}5 ; 16{,}6 ; 16{,}8 ; 16{,}9 ; 17{,}0 ; 17{,}2 ; 17{,}4 ;
17{,}5 ; 17{,}7 ; 17{,}9 ; 18{,}1 ; 18{,}4 ; 18{,}6 ; 18{,}9 ; 19{,}1}

% =========================================================
\begin{document}

% ============== PAGE DE GARDE ==============
\begin{center}
{\color{corrige}\rule{\linewidth}{1.2pt}}\\[0.3cm]
{\Huge\bfseries\color{corrige} CORRIGÉ DU CHALLENGE STATISTIQUE}\\[0.2cm]
{\Large\bfseries\color{bioaccent} STS BIOAL -- Module 5}\\[0.1cm]
{\large Statistique descriptive -- Classe de 40 échantillons sériques}\\[0.2cm]
{\color{corrige}\rule{\linewidth}{1.2pt}}
\end{center}

\vspace{0.3cm}

\begin{tcolorbox}[colback=corrigelight,colframe=corrige,title={\textbf{Mode d'emploi de ce corrigé}},fonttitle=\bfseries]
Ce document corrige l'intégralité du Challenge statistique. Il est structuré en \textbf{trois parties} :
\begin{enumerate}[leftmargin=*,itemsep=1pt]
  \item \textbf{Corrections détaillées} des 6 épreuves (calculs, justifications, pièges).
  \item \textbf{Représentations graphiques} implantées à chaque fois qu'elles éclairent le propos (histogramme, courbe des FCC, boîte à moustaches, nuage de points).
  \item \textbf{Guide de choix des graphiques} : quel graphique pour quelle situation ? Tableau comparatif et arbre de décision.
\end{enumerate}

\medskip
\textbf{Convention.} Les résultats sont donnés avec les arrondis demandés dans l'énoncé. Les étapes intermédiaires sont conservées pour la pédagogie. Les pièges sont signalés par un encadré rouge.
\end{tcolorbox}

\vspace{0.3cm}

% =========================================================
% RAPPEL DU CONTEXTE
% =========================================================
\begin{tcolorbox}[colback=bioyellow,colframe=bioaccent,title={\textbf{Rappel du contexte commun}},fonttitle=\bfseries]
Un laboratoire de bioanalyse contrôle la \textbf{concentration en glucose} (g/L) d'un lot de 40 échantillons sériques :
\vspace{0.2cm}
\centering
\begin{minipage}{0.95\linewidth}
\small\dataS
\end{minipage}
\end{tcolorbox}

\newpage

% =========================================================
% ÉPREUVE 1 -- CORRIGÉ
% =========================================================
\begin{epreuve}{CORRIGÉ 1 -- Tableau statistique \hfill \small (8 min -- 20 pts)}

\vspace{0.2cm}
\textbf{1.1. Tableau complété.}
\emph{Comptage des effectifs :} on parcourt la liste et on compte chaque valeur dans sa classe.

\vspace{0.2cm}
\renewcommand{\arraystretch}{1.45}
\begin{center}
\small
\begin{tabular}{|c|c|c|c|c|}
\hline
\rowcolor{corrige!15}
\textbf{Classe} & \textbf{Effectif $n_i$} & \textbf{Fréquence $f_i$} & \textbf{FCC} & \textbf{Centre $x_i$} \\
\hline
$[12\,;13[$ & 2 & 0{,}050 & 0{,}050 & 12{,}5 \\
\hline
$[13\,;14[$ & 4 & 0{,}100 & 0{,}150 & 13{,}5 \\
\hline
$[14\,;15[$ & 6 & 0{,}150 & 0{,}300 & 14{,}5 \\
\hline
$[15\,;16[$ & 9 & 0{,}225 & 0{,}525 & 15{,}5 \\
\hline
$[16\,;17[$ & 8 & 0{,}200 & 0{,}725 & 16{,}5 \\
\hline
$[17\,;18[$ & 6 & 0{,}150 & 0{,}875 & 17{,}5 \\
\hline
$[18\,;19[$ & 4 & 0{,}100 & 0{,}975 & 18{,}5 \\
\hline
$[19\,;20[$ & 1 & 0{,}025 & 1{,}000 & 19{,}5 \\
\hline
\rowcolor{corrige!15}
\textbf{Total} & \textbf{40} & \textbf{1{,}000} & --- & --- \\
\hline
\end{tabular}
\end{center}

\medskip
\textbf{Vérifications :} $\sum n_i = 40$ \checkmark \quad $\sum f_i = 1{,}000$ \checkmark \quad Dernière FCC $= 1{,}000$ \checkmark

\vspace{0.3cm}
\textbf{1.2. FCC de la classe $[16\,;17[$ :} $F_5 = 0{,}725$.

\medskip
\textbf{Interprétation.} \emph{72{,}5\,\% des échantillons ont une concentration en glucose strictement inférieure à 17\,g/L.}

\vspace{0.3cm}
\textbf{1.3. Proportion strictement inférieure à 15\,g/L.}

La valeur 15 est la borne supérieure de la classe $[14\,;15[$. On lit donc la FCC de cette classe :
$$F_3 = f_1+f_2+f_3 = 0{,}050+0{,}100+0{,}150 = 0{,}300.$$
\emph{Réponse :} \fbox{\textbf{30{,}0\,\%}} des échantillons ont une concentration strictement inférieure à 15\,g/L.

\vspace{0.3cm}
\textbf{1.4. Proportion supérieure ou égale à 17\,g/L.}

On somme les effectifs des classes $[17\,;18[$, $[18\,;19[$ et $[19\,;20[$ :
$$n = 6 + 4 + 1 = 11 \quad\Rightarrow\quad f = \frac{11}{40} = 0{,}275.$$
\emph{Réponse :} \fbox{\textbf{27{,}5\,\%}}.
\emph{Autre méthode :} $1 - F_5 = 1 - 0{,}725 = 0{,}275$.

\end{epreuve}

\vspace{0.3cm}

% ---------- Histogramme et FCC ----------
\begin{vignetteCours}{Représentation 1 -- Histogramme des effectifs (8 classes de largeur 1)}
\begin{center}
\begin{tikzpicture}
\begin{axis}[
  width=14cm, height=6cm,
  ybar interval, bar width=0.9,
  xlabel={Concentration en glucose (g/L)},
  ylabel={Effectif $n_i$},
  xmin=12, xmax=20, ymin=0, ymax=10,
  xtick={12,13,14,15,16,17,18,19,20},
  ytick={0,2,4,6,8,10},
  axis lines=left,
  tick label style={font=\small},
  label style={font=\small},
  every axis plot/.append style={fill=biocolor!55, draw=biocolor},
  nodes near coords, nodes near coords align={above},
  every node near coord/.append style={font=\small},
]
\addplot coordinates {(12,2)(13,4)(14,6)(15,9)(16,8)(17,6)(18,4)(19,1)(20,0)};
\end{axis}
\end{tikzpicture}
\end{center}
\textbf{Intérêt pédagogique.} L'histogramme \textbf{montre la forme de la distribution} :
\begin{itemize}[leftmargin=*,itemsep=1pt]
  \item \textbf{mode} visible à $[15\,;16[$ (hauteur maximale) ;
  \item distribution \textbf{quasi symétrique} autour de la classe modale ;
  \item \textbf{pas d'outlier} évident (la classe $[19\,;20[$ n'a qu'un seul individu mais reste dans la continuité).
\end{itemize}
\emph{C'est le graphique de référence pour une variable continue regroupée en classes.}
\end{vignetteCours}

\vspace{0.3cm}

\begin{vignetteCours}{Représentation 2 -- Courbe des fréquences cumulées croissantes (FCC)}
\begin{center}
\begin{tikzpicture}
\begin{axis}[
  width=14cm, height=6.5cm,
  xlabel={Concentration en glucose (g/L)},
  ylabel={Fréquence cumulée croissante},
  xmin=12, xmax=20, ymin=0, ymax=1.05,
  xtick={12,13,14,15,16,17,18,19,20},
  ytick={0,0.25,0.5,0.75,1},
  yticklabels={0, 0{,}25, 0{,}50, 0{,}75, 1},
  grid=both, grid style={gray!25},
  axis lines=left,
  tick label style={font=\small},
  label style={font=\small},
]
\addplot[thick, color=biocolor, mark=*, mark size=2.5pt]
  coordinates {(12,0)(13,0.050)(14,0.150)(15,0.300)(16,0.525)(17,0.725)(18,0.875)(19,0.975)(20,1.000)};
\draw[dashed, bioaccent, thick] (axis cs:12,0.25) -- (axis cs:14.667,0.25) -- (axis cs:14.667,0);
\draw[dashed, bioaccent, thick] (axis cs:12,0.50) -- (axis cs:15.889,0.50) -- (axis cs:15.889,0);
\draw[dashed, bioaccent, thick] (axis cs:12,0.75) -- (axis cs:17.167,0.75) -- (axis cs:17.167,0);
\node[font=\scriptsize, bioaccent] at (axis cs:14.3,0.30) {$Q_1\!\approx\!14{,}67$};
\node[font=\scriptsize, bioaccent] at (axis cs:15.5,0.57) {$M_e\!\approx\!15{,}89$};
\node[font=\scriptsize, bioaccent] at (axis cs:16.6,0.82) {$Q_3\!\approx\!17{,}17$};
\end{axis}
\end{tikzpicture}
\end{center}
\textbf{Intérêt pédagogique.} La courbe des FCC permet de \textbf{lire directement} $Q_1$, $M_e$, $Q_3$ par interpolation linéaire, et de répondre à toute question du type « quel pourcentage a une valeur $< x$ ? ». C'est le \textbf{complément naturel de l'histogramme}.
\end{vignetteCours}

\newpage

% =========================================================
% ÉPREUVE 2 -- CORRIGÉ
% =========================================================
\begin{epreuve}{CORRIGÉ 2 -- Représentations graphiques \hfill \small (8 min -- 15 pts)}

\vspace{0.2cm}
\textbf{2.1. et 2.2.} Voir les deux graphiques de la page précédente (histogramme + courbe des FCC).

\vspace{0.2cm}
\textbf{2.3. Détermination des cinq indicateurs.}

\begin{center}
\renewcommand{\arraystretch}{1.35}
\begin{tabular}{|l|c|l|}
\hline
\rowcolor{corrige!15}
\textbf{Indicateur} & \textbf{Valeur} & \textbf{Justification} \\
\hline
Minimum & 12{,}1 & 1\textsuperscript{re} valeur de la série ordonnée \\
\hline
Maximum & 19{,}1 & Dernière valeur de la série ordonnée \\
\hline
Médiane $M_e$ & 15{,}85 & Moyenne des rangs 20 et 21 : $(15{,}8+15{,}9)/2$ \\
\hline
Premier quartile $Q_1$ & 14{,}5 & Rang $\lceil 40/4 \rceil = 10$ \\
\hline
Troisième quartile $Q_3$ & 17{,}0 & Rang $\lceil 3\times40/4 \rceil = 30$ \\
\hline
\end{tabular}
\end{center}

\medskip
\textbf{2.4. Boîte à moustaches.} Placement des 5 valeurs : $\min=12{,}1$ ; $Q_1=14{,}5$ ; $M_e=15{,}85$ ; $Q_3=17{,}0$ ; $\max=19{,}1$.

\begin{center}
\begin{tikzpicture}[xscale=1.8, yscale=0.9]
  \draw[thick,<->] (11.8,0) -- (19.6,0);
  \foreach \x in {12,13,14,15,16,17,18,19}
    \draw[thick] (\x,0.12) -- (\x,-0.12) node[below,font=\small]{\x};
  % Boîte Q1 - Q3
  \draw[thick, fill=corrige!30] (14.5,-0.5) rectangle (17.0,0.5);
  % Médiane
  \draw[thick, corrige] (15.85,-0.5) -- (15.85,0.5);
  % Moustaches
  \draw[thick] (12.1,0) -- (14.5,0);
  \draw[thick] (17.0,0) -- (19.1,0);
  \draw[thick] (12.1,-0.15) -- (12.1,0.15);
  \draw[thick] (19.1,-0.15) -- (19.1,0.15);
  % Étiquettes
  \node[above, font=\scriptsize] at (12.1,0.2) {min=12{,}1};
  \node[above, font=\scriptsize] at (14.5,0.6) {$Q_1$=14{,}5};
  \node[above, font=\scriptsize, corrige] at (15.85,0.6) {$M_e$=15{,}85};
  \node[above, font=\scriptsize] at (17.0,0.6) {$Q_3$=17{,}0};
  \node[above, font=\scriptsize] at (19.1,0.2) {max=19{,}1};
  \node[below, font=\scriptsize] at (15.85,-0.75) {g/L};
\end{tikzpicture}
\end{center}

\medskip
\textbf{2.5. Symétrie.} On compare les deux demi-boîtes et les deux moustaches :
\begin{itemize}[leftmargin=*,itemsep=1pt]
  \item $M_e - Q_1 = 15{,}85 - 14{,}5 = 1{,}35$ \quad et \quad $Q_3 - M_e = 17{,}0 - 15{,}85 = 1{,}15$ : \textbf{écarts proches} ;
  \item moustache gauche : $Q_1 - \min = 14{,}5 - 12{,}1 = 2{,}4$ \quad et \quad moustache droite : $\max - Q_3 = 19{,}1 - 17{,}0 = 2{,}1$ : \textbf{longueurs proches}.
\end{itemize}
\emph{Conclusion.} La distribution est \textbf{approximativement symétrique} (léger étalement vers la gauche, mais négligeable).

\end{epreuve}

\newpage

% =========================================================
% ÉPREUVE 3 -- CORRIGÉ
% =========================================================
\begin{epreuve}{CORRIGÉ 3 -- Indicateurs de position \hfill \small (8 min -- 20 pts)}

\vspace{0.2cm}
\textbf{3.1. Moyenne.} On calcule d'abord $\sum n_i x_i$ à partir des centres de classes.

\begin{center}
\small
\renewcommand{\arraystretch}{1.35}
\begin{tabular}{|c|c|c|c|}
\hline
\rowcolor{corrige!15}
\textbf{Classe} & $x_i$ & $n_i$ & $n_i x_i$ \\
\hline
$[12\,;13[$ & 12{,}5 & 2 & 25{,}0 \\
\hline
$[13\,;14[$ & 13{,}5 & 4 & 54{,}0 \\
\hline
$[14\,;15[$ & 14{,}5 & 6 & 87{,}0 \\
\hline
$[15\,;16[$ & 15{,}5 & 9 & 139{,}5 \\
\hline
$[16\,;17[$ & 16{,}5 & 8 & 132{,}0 \\
\hline
$[17\,;18[$ & 17{,}5 & 6 & 105{,}0 \\
\hline
$[18\,;19[$ & 18{,}5 & 4 & 74{,}0 \\
\hline
$[19\,;20[$ & 19{,}5 & 1 & 19{,}5 \\
\hline
\rowcolor{corrige!15}
\textbf{Total} & --- & \textbf{40} & $\sum n_i x_i = \mathbf{636{,}0}$ \\
\hline
\end{tabular}
\end{center}

$$\boxed{\bar{x} = \frac{\sum n_i x_i}{N} = \frac{636{,}0}{40} = 15{,}90\ \text{g/L}}$$

\vspace{0.3cm}
\textbf{3.2. Médiane.} On utilise les rangs. $N = 40$ est \emph{pair}, donc :
$$M_e = \frac{\text{valeur de rang } 20 + \text{valeur de rang } 21}{2}.$$

La série ordonnée comporte (cumul) : rangs 1--2 : $[12\,;13[$ ; rangs 3--6 : $[13\,;14[$ ; rangs 7--12 : $[14\,;15[$ ; rangs 13--21 : $[15\,;16[$.

Les rangs 20 et 21 tombent tous deux dans la classe $[15\,;16[$ : rang 20 = $15{,}8$, rang 21 = $15{,}9$.
$$\boxed{M_e = \frac{15{,}8+15{,}9}{2} = 15{,}85\ \text{g/L}}$$

\vspace{0.3cm}
\textbf{3.3. Quartiles.}
\begin{itemize}[leftmargin=*,itemsep=1pt]
  \item $Q_1$ : rang $\lceil N/4 \rceil = 10$ $\Rightarrow$ \fbox{$Q_1 = 14{,}5$ g/L}. \\
  \emph{Signification concrète :} 25\,\% des échantillons ont une concentration inférieure ou égale à 14{,}5 g/L (et 75\,\% supérieure ou égale).
  \item $Q_3$ : rang $\lceil 3N/4 \rceil = 30$ $\Rightarrow$ \fbox{$Q_3 = 17{,}0$ g/L}.
\end{itemize}

\vspace{0.3cm}
\textbf{3.4. Mode et classe modale.}
\begin{itemize}[leftmargin=*,itemsep=1pt]
  \item Il s'agit d'une série \textbf{continue regroupée en classes} : le mode est une \textbf{classe modale}.
  \item La classe de plus grand effectif est $[15\,;16[$ avec $n = 9$.
  \item \fbox{Classe modale : $[15\,;16[$}. Il n'y a pas de mode unique (valeur ponctuelle), mais une classe modale unique.
\end{itemize}

\vspace{0.3cm}
\textbf{3.5. Vrai/Faux.}
\begin{itemize}[leftmargin=*,itemsep=2pt]
  \item \emph{« Comme la moyenne et la médiane sont proches, la distribution est approximativement symétrique. »} \\
  \textbf{VRAI.} $\bar{x} = 15{,}90$ et $M_e = 15{,}85$ : l'écart est de $0{,}05$, négligeable. Ceci confirme la lecture de la boîte à moustaches (question 2.5). \emph{Attention :} une proximité moyenne/médiane est une \emph{condition nécessaire} mais non suffisante ; elle doit être confirmée par l'histogramme.
  \item \emph{« 50\,\% des échantillons ont une concentration comprise entre $Q_1$ et $Q_3$. »} \\
  \textbf{VRAI.} Par définition, $Q_1$ et $Q_3$ encadrent les 50\,\% centraux de la population. Ici : $14{,}5 \le x \le 17{,}0$ contient 50\,\% des valeurs.
\end{itemize}

\end{epreuve}

\newpage

% =========================================================
% ÉPREUVE 4 -- CORRIGÉ
% =========================================================
\begin{epreuve}{CORRIGÉ 4 -- Indicateurs de dispersion \hfill \small (7 min -- 15 pts)}

\vspace{0.2cm}
\textbf{4.1. Étendue.}
$$\boxed{E = \max - \min = 19{,}1 - 12{,}1 = 7{,}0\ \text{g/L}}$$

\vspace{0.3cm}
\textbf{4.2. Écart interquartile.}
$$\boxed{I_Q = Q_3 - Q_1 = 17{,}0 - 14{,}5 = 2{,}5\ \text{g/L}}$$

\vspace{0.3cm}
\textbf{4.3. Variance et écart-type.} On utilise la formule rapide $V = \dfrac{1}{N}\sum n_i x_i^2 - \bar{x}^2$.

\begin{center}
\small
\renewcommand{\arraystretch}{1.35}
\begin{tabular}{|c|c|c|c|c|}
\hline
\rowcolor{corrige!15}
\textbf{Classe} & $x_i$ & $n_i$ & $x_i^2$ & $n_i x_i^2$ \\
\hline
$[12\,;13[$ & 12{,}5 & 2 & 156{,}25 & 312{,}50 \\
\hline
$[13\,;14[$ & 13{,}5 & 4 & 182{,}25 & 729{,}00 \\
\hline
$[14\,;15[$ & 14{,}5 & 6 & 210{,}25 & 1\,261{,}50 \\
\hline
$[15\,;16[$ & 15{,}5 & 9 & 240{,}25 & 2\,162{,}25 \\
\hline
$[16\,;17[$ & 16{,}5 & 8 & 272{,}25 & 2\,178{,}00 \\
\hline
$[17\,;18[$ & 17{,}5 & 6 & 306{,}25 & 1\,837{,}50 \\
\hline
$[18\,;19[$ & 18{,}5 & 4 & 342{,}25 & 1\,369{,}00 \\
\hline
$[19\,;20[$ & 19{,}5 & 1 & 380{,}25 & 380{,}25 \\
\hline
\rowcolor{corrige!15}
\textbf{Total} & --- & \textbf{40} & --- & $\mathbf{10\,230{,}0}$ \\
\hline
\end{tabular}
\end{center}

$$V = \frac{10\,230{,}0}{40} - (15{,}90)^2 = 255{,}75 - 252{,}81 = 2{,}94$$
$$\boxed{V \approx 2{,}94\ (\text{g/L})^2} \qquad \boxed{\sigma = \sqrt{2{,}94} \approx 1{,}715\ \text{g/L}}$$

\vspace{0.3cm}
\textbf{4.4. Interprétation de l'écart-type.}

\emph{« En moyenne, les concentrations en glucose des 40 échantillons s'écartent d'environ $1{,}7$ g/L de la concentration moyenne (15{,}90 g/L). »}

\medskip
Cette valeur traduit l'homogénéité du lot : un écart-type faible par rapport à la moyenne ($\sigma/\bar{x} \approx 10{,}8\,\%$) indique un lot \emph{relativement homogène}, sans grande variabilité.

\end{epreuve}

\newpage

% =========================================================
% ÉPREUVE 5 -- CORRIGÉ
% =========================================================
\begin{epreuve}{CORRIGÉ 5 -- Ajustement affine et corrélation \hfill \small (12 min -- 20 pts)}

\vspace{0.2cm}
\textbf{5.1. Moyennes.}
$$\bar{x} = \frac{0+2+4+6+8+10}{6} = \frac{30}{6} = 5{,}00$$
$$\bar{y} = \frac{0{,}05+0{,}22+0{,}38+0{,}57+0{,}71+0{,}90}{6} = \frac{2{,}83}{6} \approx 0{,}4717$$

\vspace{0.3cm}
\textbf{5.2. Sommes.} On construit un tableau de calcul.

\begin{center}
\small
\renewcommand{\arraystretch}{1.35}
\begin{tabular}{|c|c|c|c|c|c|}
\hline
\rowcolor{corrige!15}
$i$ & $x_i$ & $y_i$ & $x_i^2$ & $y_i^2$ & $x_i y_i$ \\
\hline
1 & 0 & 0{,}05 & 0 & 0{,}0025 & 0 \\
\hline
2 & 2 & 0{,}22 & 4 & 0{,}0484 & 0{,}44 \\
\hline
3 & 4 & 0{,}38 & 16 & 0{,}1444 & 1{,}52 \\
\hline
4 & 6 & 0{,}57 & 36 & 0{,}3249 & 3{,}42 \\
\hline
5 & 8 & 0{,}71 & 64 & 0{,}5041 & 5{,}68 \\
\hline
6 & 10 & 0{,}90 & 100 & 0{,}8100 & 9{,}00 \\
\hline
\rowcolor{corrige!15}
\textbf{Total} & $\sum x_i=30$ & $\sum y_i=2{,}83$ & $\mathbf{220}$ & $\mathbf{1{,}8343}$ & $\mathbf{20{,}06}$ \\
\hline
\end{tabular}
\end{center}

\vspace{0.3cm}
\textbf{5.3. Droite de régression $y = ax + b$.}

\emph{Covariance :}
$$\mathrm{Cov}(x,y) = \frac{\sum x_i y_i}{N} - \bar{x}\,\bar{y} = \frac{20{,}06}{6} - 5 \times 0{,}4717 = 3{,}3433 - 2{,}3585 = 0{,}9848$$

\emph{Variances :}
$$V(x) = \frac{220}{6} - 5^2 = 36{,}6667 - 25 = 11{,}6667$$
$$V(y) = \frac{1{,}8343}{6} - (0{,}4717)^2 = 0{,}3057 - 0{,}2225 = 0{,}0832$$

\emph{Coefficients :}
$$a = \frac{\mathrm{Cov}(x,y)}{V(x)} = \frac{0{,}9848}{11{,}6667} \approx 0{,}0844$$
$$b = \bar{y} - a\bar{x} = 0{,}4717 - 0{,}0844 \times 5 = 0{,}4717 - 0{,}4220 \approx 0{,}050$$

$$\boxed{y = 0{,}0844\,x + 0{,}050}$$

\vspace{0.3cm}
\textbf{5.4. Coefficient de corrélation linéaire.}
$$r = \frac{\mathrm{Cov}(x,y)}{\sqrt{V(x)\,V(y)}} = \frac{0{,}9848}{\sqrt{11{,}6667 \times 0{,}0832}} = \frac{0{,}9848}{0{,}9855} \approx 0{,}999$$

$$\boxed{r \approx 0{,}999}$$

\vspace{0.3cm}
\textbf{5.5. Adéquation du modèle linéaire.}

$|r| = 0{,}999$ est \textbf{très proche de 1} : la liaison linéaire entre la concentration et l'absorbance est \textbf{quasi parfaite}. Le modèle linéaire est donc \textbf{parfaitement adapté}.

\medskip
\emph{Vérification graphique} (voir nuage de points ci-dessous) : les 6 points sont presque exactement alignés.

\vspace{0.3cm}
\textbf{5.6. Application.} Pour $y = 0{,}80$ :
$$0{,}80 = 0{,}0844\,x + 0{,}050 \quad\Rightarrow\quad x = \frac{0{,}80 - 0{,}050}{0{,}0844} \approx 8{,}89$$

$$\boxed{x \approx 8{,}89\ \text{g/L}}$$

\end{epreuve}

\vspace{0.3cm}

\begin{vignetteCours}{Représentation 3 -- Nuage de points et droite de régression $y = 0{,}0844\,x + 0{,}050$}
\begin{center}
\begin{tikzpicture}
\begin{axis}[
  width=14cm, height=7cm,
  xlabel={Concentration $x$ (g/L)},
  ylabel={Absorbance $y$},
  xmin=0, xmax=11, ymin=0, ymax=1.05,
  xtick={0,2,4,6,8,10},
  ytick={0,0.2,0.4,0.6,0.8,1.0},
  grid=both, grid style={gray!25},
  axis lines=left,
  tick label style={font=\small},
  label style={font=\small},
]
\addplot[only marks, mark=*, mark size=3.2pt, color=bioaccent]
  coordinates {(0,0.05)(2,0.22)(4,0.38)(6,0.57)(8,0.71)(10,0.90)};
\addplot[thick, color=corrige, domain=0:10.5, samples=2] {0.0844*x + 0.050};
\node[font=\small, corrige] at (axis cs:7.8,0.35) {$y = 0{,}0844\,x + 0{,}050$};
\node[font=\small, corrige] at (axis cs:7.8,0.20) {$r \approx 0{,}999$};
\draw[dashed, bioaccent] (axis cs:8.89,0) -- (axis cs:8.89,0.80) -- (axis cs:0,0.80);
\node[font=\scriptsize, bioaccent] at (axis cs:9.5,0.68) {$(8{,}89\,;\,0{,}80)$};
\end{axis}
\end{tikzpicture}
\end{center}
\textbf{Intérêt pédagogique.} Le nuage de points est \textbf{le graphique indispensable} pour :
\begin{itemize}[leftmargin=*,itemsep=1pt]
  \item \textbf{vérifier la linéarité} avant tout calcul de régression ;
  \item \textbf{repérer les outliers} (points très éloignés de la tendance) ;
  \item \textbf{visualiser l'étalonnage} et lire une valeur inconnue (ici, la flèche en pointillés montre la lecture pour $y=0{,}80$).
\end{itemize}
\end{vignetteCours}

\newpage

% =========================================================
% ÉPREUVE 6 -- CORRIGÉ
% =========================================================
\begin{epreuve}{CORRIGÉ 6 -- Utilisation du tableur \hfill \small (7 min -- 10 pts)}

\vspace{0.2cm}
\textbf{Formules exactes à saisir} (plage \texttt{A2:A41}) :

\vspace{0.2cm}
\begin{center}
\renewcommand{\arraystretch}{1.5}
\begin{tabular}{|l|l|l|}
\hline
\rowcolor{corrige!15}
\textbf{Question} & \textbf{Fonction} & \textbf{Formule exacte} \\
\hline
a. Effectif total & NB & \texttt{=NB(A2:A41)} \\
\hline
b. Moyenne & MOYENNE & \texttt{=MOYENNE(A2:A41)} \\
\hline
c. Médiane & MEDIANE & \texttt{=MEDIANE(A2:A41)} \\
\hline
d. Premier quartile $Q_1$ & QUARTILE & \texttt{=QUARTILE(A2:A41;1)} \\
\hline
e. Troisième quartile $Q_3$ & QUARTILE & \texttt{=QUARTILE(A2:A41;3)} \\
\hline
f. Écart-type & ECARTYPE & \texttt{=ECARTYPE(A2:A41)} \\
\hline
g. Minimum & MIN & \texttt{=MIN(A2:A41)} \\
\hline
h. Maximum & MAX & \texttt{=MAX(A2:A41)} \\
\hline
\end{tabular}
\end{center}

\medskip
\emph{Remarque.} Selon la version d'Excel, on peut aussi écrire \texttt{=QUARTILE.INCLURE(A2:A41;1)} (qui donne le même résultat que \texttt{QUARTILE}). Pour l'écart-type d'une \emph{population}, utiliser \texttt{=ECARTYPEP(A2:A41)}.

\vspace{0.3cm}
\textbf{i. Deux méthodes pour tracer une boîte à moustaches.}

\begin{enumerate}[leftmargin=*,itemsep=2pt]
  \item \textbf{Méthode manuelle} : calculer séparément $\min$, $Q_1$, $M_e$, $Q_3$, $\max$ (via les fonctions précédentes), puis construire le graphique à la main ou avec l'outil « graphique en chandeliers » (\emph{stock chart}).
  \item \textbf{Méthode automatique} : sélectionner la plage \texttt{A2:A41}, puis \texttt{Insertion > Graphiques > Boîte à moustaches} (Excel 2016 et versions ultérieures). Dans LibreOffice Calc : \texttt{Insertion > Graphique > Boîte à moustaches}.
\end{enumerate}

\vspace{0.3cm}
\textbf{j. Recopie vers le bas.}
\begin{itemize}[leftmargin=*,itemsep=1pt]
  \item En \texttt{C1}, on a saisi \texttt{=MOYENNE(A2:A41)}.
  \item En recopiant vers \texttt{C2}, on obtient \texttt{=MOYENNE(A3:A42)}.
  \item \textbf{Le résultat n'est pas pertinent} : la plage se décale automatiquement (référence \emph{relative}), et les cellules \texttt{A3:A42} ne correspondent plus à la série (perte de la première valeur, prise en compte d'une cellule vide ou parasite).
  \item \textbf{Solution} : utiliser des \textbf{références absolues} en ancrant les bornes avec le symbole \texttt{\$} :
  $$\boxed{\texttt{=MOYENNE(\$A\$2:\$A\$41)}}$$
  Ainsi, la recopie en \texttt{C2}, \texttt{C3}... conserve la plage correcte.
\end{itemize}

\end{epreuve}

\newpage

% =========================================================
% GUIDE DE CHOIX DES GRAPHIQUES
% =========================================================
\section*{\color{corrige}Guide de choix des représentations graphiques}
\emph{Quel graphique pour quelle situation ? Synthèse pédagogique.}

\vspace{0.3cm}

\begin{tcolorbox}[enhanced, colback=corrigelight, colframe=corrige,
                  title={\faChartBar\ \ \textbf{Tableau comparatif des graphiques}},
                  coltitle=white, fonttitle=\bfseries\large,
                  boxrule=0.8pt, arc=4pt, left=8pt, right=8pt, top=8pt, bottom=8pt,
                  drop fuzzy shadow=black!15]
\small
\renewcommand{\arraystretch}{1.4}
\begin{tabular}{|p{2.3cm}|p{2.6cm}|p{2.9cm}|p{3.3cm}|p{2.5cm}|}
\hline
\rowcolor{corrige!25}
\textbf{Graphique} & \textbf{Type de variable} & \textbf{Ce qu'il montre} & \textbf{Intérêt principal} & \textbf{Limites} \\
\hline
\textbf{Bâtons} & Qualitative ou quantitative discrète & Effectif / fréquence de chaque modalité & Compare les catégories ; repère le mode & Inadapté aux variables continues \\
\hline
\textbf{Circulaire} & Qualitative (peu de modalités) & Part de chaque catégorie dans le total & Vision « composition » parlante & Illisible au-delà de 6--7 modalités \\
\hline
\textbf{Histogramme} & Quantitative continue (classes) & Forme de la distribution (symétrie, mode, outliers) & Révèle la structure globale ; oriente le choix moyenne/médiane & Sensible au choix des classes \\
\hline
\textbf{Courbe des FCC} & Quantitative (discrète ou continue) & Cumul des fréquences ; lecture directe de $M_e$, $Q_1$, $Q_3$ & Répond à « quel \% est $< x$ ? » ; complément de l'histogramme & Moins visuelle seule \\
\hline
\textbf{Boîte à moustaches} & Quantitative (une ou plusieurs séries) & 5 indicateurs : min, $Q_1$, $M_e$, $Q_3$, max ; outliers & Compare plusieurs séries en un coup d'œil ; robuste & Ne montre pas la forme exacte (bimodalité invisible) \\
\hline
\textbf{Nuage de points + droite} & Deux variables quantitatives & Liaison entre deux variables ; tendance linéaire & Prédiction, étalonnage ; détection d'outliers et de non-linéarité & $r$ ne mesure que la linéarité \\
\hline
\end{tabular}
\end{tcolorbox}

\vspace{0.3cm}

\begin{vignetteMethode}{Arbre de décision -- Quel graphique choisir ?}
\begin{enumerate}[leftmargin=*,itemsep=2pt]
  \item \textbf{Une seule variable, qualitative ou discrète ?} \textrightarrow\ Diagramme en bâtons (ou circulaire si peu de modalités et objectif de composition).
  \item \textbf{Une seule variable, quantitative continue ?}
    \begin{itemize}[leftmargin=*,itemsep=1pt]
      \item Pour voir la \emph{forme} de la distribution \textrightarrow\ \textbf{Histogramme}.
      \item Pour lire $M_e$, $Q_1$, $Q_3$ sans calcul \textrightarrow\ \textbf{Courbe des FCC}.
      \item Pour \emph{comparer plusieurs séries} \textrightarrow\ \textbf{Boîtes à moustaches juxtaposées}.
    \end{itemize}
  \item \textbf{Deux variables quantitatives ?} \textrightarrow\ \textbf{Nuage de points}, puis droite de régression si la tendance est linéaire.
  \item \textbf{Règle d'or} : \emph{toujours tracer d'abord, calculer ensuite}. Un graphique mal choisi peut masquer une information essentielle.
\end{enumerate}
\end{vignetteMethode}

\vspace{0.3cm}

\begin{vignetteMethode}{Quel graphique pour notre sujet ? Réponses concrètes}
\begin{itemize}[leftmargin=*,itemsep=2pt]
  \item \textbf{Pour montrer la répartition des 40 concentrations} \textrightarrow\ \textbf{histogramme} (Réprésentation 1).
  \item \textbf{Pour lire directement $M_e$, $Q_1$, $Q_3$} \textrightarrow\ \textbf{courbe des FCC} (Représentation 2).
  \item \textbf{Pour visualiser la dispersion et la symétrie} \textrightarrow\ \textbf{boîte à moustaches} (voir question 2.4).
  \item \textbf{Pour l'étalonnage du spectrophotomètre} \textrightarrow\ \textbf{nuage de points} + droite de régression (Représentation 3).
  \item \textbf{Pour comparer plusieurs lots} (non demandé ici, mais utile) \textrightarrow\ \textbf{boîtes à moustaches juxtaposées}.
\end{itemize}
\end{vignetteMethode}

\vspace{0.3cm}

\begin{vignetteAlerte}{Pièges classiques à éviter}
\begin{itemize}[leftmargin=*,itemsep=2pt]
  \item \textbf{Histogramme vs diagramme en bâtons} : l'histogramme a des barres \emph{jointives} (variable continue) ; le diagramme en bâtons a des barres \emph{séparées} (variable discrète ou qualitative).
  \item \textbf{Aires proportionnelles} : dans un histogramme, ce sont les \emph{aires} qui sont proportionnelles aux effectifs, pas les hauteurs (sauf si les classes ont même amplitude).
  \item \textbf{Courbe des FCC} : elle est \emph{toujours croissante}, de 0 à 1. Une courbe qui descend signale une erreur de calcul.
  \item \textbf{Boîte à moustaches} : elle \emph{ne montre pas} la forme exacte de la distribution. Deux séries très différentes (unimodale vs bimodale) peuvent avoir la même boîte.
  \item \textbf{Nuage de points} : toujours le tracer \emph{avant} de calculer $r$. Un $r \approx 0$ n'exclut pas une relation non linéaire (parabole, exponentielle...).
  \item \textbf{$r$ proche de 1} : ne prouve pas une relation de \emph{causalité}, seulement une corrélation \emph{linéaire} forte.
\end{itemize}
\end{vignetteAlerte}

\newpage

% =========================================================
% RÉCAPITULATIF DES RÉSULTATS
% =========================================================
\section*{\color{corrige}Récapitulatif des résultats numériques}

\vspace{0.3cm}

\begin{tcolorbox}[enhanced, colback=corrigelight, colframe=corrige,
                  title={\faListOl\ \ \textbf{Tous les résultats en un coup d'œil}},
                  coltitle=white, fonttitle=\bfseries\large,
                  boxrule=0.8pt, arc=4pt, left=8pt, right=8pt, top=8pt, bottom=8pt,
                  drop fuzzy shadow=black!15]
\small
\renewcommand{\arraystretch}{1.4}
\begin{tabular}{|l|l|l|}
\hline
\rowcolor{corrige!25}
\textbf{Épreuve} & \textbf{Quantité} & \textbf{Résultat} \\
\hline
\multirow{4}{*}{1} & Effectifs $n_i$ & 2 ; 4 ; 6 ; 9 ; 8 ; 6 ; 4 ; 1 \\
\cline{2-3}
 & FCC de $[16\,;17[$ & 0{,}725 \\
\cline{2-3}
 & Proportion $<15$ g/L & 30{,}0\,\% \\
\cline{2-3}
 & Proportion $\ge 17$ g/L & 27{,}5\,\% \\
\hline
\multirow{5}{*}{2 et 3} & Minimum & 12{,}1 g/L \\
\cline{2-3}
 & Maximum & 19{,}1 g/L \\
\cline{2-3}
 & Médiane $M_e$ & 15{,}85 g/L \\
\cline{2-3}
 & Premier quartile $Q_1$ & 14{,}5 g/L \\
\cline{2-3}
 & Troisième quartile $Q_3$ & 17{,}0 g/L \\
\hline
\multirow{3}{*}{3} & Moyenne $\bar{x}$ & 15{,}90 g/L \\
\cline{2-3}
 & Mode & Classe $[15\,;16[$ \\
\cline{2-3}
 & Symétrie & approximativement symétrique \\
\hline
\multirow{3}{*}{4} & Étendue & 7{,}0 g/L \\
\cline{2-3}
 & Écart interquartile $I_Q$ & 2{,}5 g/L \\
\cline{2-3}
 & Variance $V$ / écart-type $\sigma$ & 2{,}94 / 1{,}715 g/L \\
\hline
\multirow{4}{*}{5} & Moyennes & $\bar{x}=5{,}00$ ; $\bar{y}=0{,}4717$ \\
\cline{2-3}
 & Pente $a$ / ordonnée $b$ & 0{,}0844 / 0{,}050 \\
\cline{2-3}
 & Coefficient $r$ & 0{,}999 \\
\cline{2-3}
 & Estimation pour $y=0{,}80$ & $x \approx 8{,}89$ g/L \\
\hline
6 & Point critique & Références absolues \texttt{\$A\$2:\$A\$41} \\
\hline
\end{tabular}
\end{tcolorbox}

\vspace{0.5cm}

\begin{tcolorbox}[colback=bioyellow!50,colframe=bioaccent,title={\textbf{Pour conclure : deux leçons méthodologiques}},fonttitle=\bfseries]
\small
\textbf{Leçon 1.} \emph{Toujours tracer les données avant de calculer.} Histogramme, boîte à moustaches et nuage de points sont vos trois alliés. Ce sont eux qui révèlent la forme de la distribution, les outliers et la linéarité — toutes choses qu'aucune formule ne peut remplacer.

\medskip
\textbf{Leçon 2.} \emph{Choisir le bon indicateur selon la situation.} Sur une distribution symétrique sans outlier (cas de ce sujet), moyenne et médiane racontent la même histoire. Sur une distribution asymétrique ou avec outliers, seuls la médiane et l'écart interquartile sont fiables.

\medskip
\emph{« Un graphique vaut mille tableaux ; un bon graphique vaut mille calculs. »}
\end{tcolorbox}

\end{document}